\begin{example}
	Based on the operations of addition, subtraction, multiplication, division, and exponentiation defined above, we present an example of an uncertain function:

	$$\hbd{\exp}{X} := \sum_{n=0}^{\infty} \left(\frac{X^n}{n!}\right) = 1 + X + \frac{X^2}{2!} + \frac{X^3}{3!} + \cdots
	$$
	And we compute a specific example $$\hbd{\exp}{\s{0,1}} \subset [e^0,e^1]$$
	For instance, with $n=8$ and accuracy up to 5 decimal places, we compute $\exp_m\left(\{0,1\}\right)$ approximately equal to
	$\{$1.00000, 1.00002, 1.00020, 1.00022, 1.00139, 1.00141, 1.00159, 1.00161, 1.00833, 1.00836, 1.00853, 1.00856, 1.00972, 1.00975, 1.00992, 1.00995, 1.04167, 1.04169, 1.04187, 1.04189, 1.04306, 1.04308, 1.04325, 1.04328, 1.05000, 1.05002, 1.05020, 1.05022, 1.05139, 1.05141, 1.05159, 1.05161, 1.16667, 1.16669, 1.16687, 1.16689, 1.16806, 1.16808, 1.16825, 1.16828, 1.17500, 1.17502, 1.17520, 1.17522, 1.17639, 1.17641, 1.17659, 1.17661, 1.20833, 1.20836, 1.20853, 1.20856, 1.20972, 1.20975, 1.20992, 1.20995, 1.21667, 1.21669, 1.21687, 1.21689, 1.21806, 1.21808, 1.21825, 1.21828, 1.50000, 1.50002, 1.50020, 1.50022, 1.50139, 1.50141, 1.50159, 1.50161, 1.50833, 1.50836, 1.50853, 1.50856, 1.50972, 1.50975, 1.50992, 1.50995, 1.54167, 1.54169, 1.54187, 1.54189, 1.54306, 1.54308, 1.54325, 1.54328, 1.55000, 1.55002, 1.55020, 1.55022, 1.55139, 1.55141, 1.55159, 1.55161, 1.66667, 1.66669, 1.66687, 1.66689, 1.66806, 1.66808, 1.66825, 1.66828, 1.67500, 1.67502, 1.67520, 1.67522, 1.67639, 1.67641, 1.67659, 1.67661, 1.70833, 1.70836, 1.70853, 1.70856, 1.70972, 1.70975, 1.70992, 1.70995, 1.71667, 1.71669, 1.71687, 1.71689, 1.71806, 1.71808, 1.71825, 1.71828, 2.00000, 2.00002, 2.00020, 2.00022, 2.00139, 2.00141, 2.00159, 2.00161, 2.00833, 2.00836, 2.00853, 2.00856, 2.00972, 2.00975, 2.00992, 2.00995, 2.04167, 2.04169, 2.04187, 2.04189, 2.04306, 2.04308, 2.04325, 2.04328, 2.05000, 2.05002, 2.05020, 2.05022, 2.05139, 2.05141, 2.05159, 2.05161, 2.16667, 2.16669, 2.16687, 2.16689, 2.16806, 2.16808, 2.16825, 2.16828, 2.17500, 2.17502, 2.17520, 2.17522, 2.17639, 2.17641, 2.17659, 2.17661, 2.20833, 2.20836, 2.20853, 2.20856, 2.20972, 2.20975, 2.20992, 2.20995, 2.21667, 2.21669, 2.21687, 2.21689, 2.21806, 2.21808, 2.21825, 2.21828, 2.50000, 2.50002, 2.50020, 2.50022, 2.50139, 2.50141, 2.50159, 2.50161, 2.50833, 2.50836, 2.50853, 2.50856, 2.50972, 2.50975, 2.50992, 2.50995, 2.54167, 2.54169, 2.54187, 2.54189, 2.54306, 2.54308, 2.54325, 2.54328, 2.55000, 2.55002, 2.55020, 2.55022, 2.55139, 2.55141, 2.55159, 2.55161, 2.66667, 2.66669, 2.66687, 2.66689, 2.66806, 2.66808, 2.66825, 2.66828, 2.67500, 2.67502, 2.67520, 2.67522, 2.67639, 2.67641, 2.67659, 2.67661, 2.70833, 2.70836, 2.70853, 2.70856, 2.70972, 2.70975, 2.70992, 2.70995, 2.71667, 2.71669, 2.71687, 2.71689, 2.71806, 2.71808, 2.71825, 2.71828$\}_u$

	\noindent Containing $2^8 = 256$ elements.
\end{example}
\begin{example}
	Returning to the example $\exp_m\s{0,1} \subseteq [e^0,e]$, it is easily seen that
	$$\Tr(e \in \exp_m\s{0,1}) = \left\{\sum_{n=0}^{\infty}\frac{1}{n!}\right\} $$
	Other transcendental numbers lying in the interval $[1,e]$ also possess a unique trace, which is the formal formula of that number.
\end{example}
\begin{example}
	$$\hbd{c}{X} = \sum_{k=1}^{\infty} \frac{X}{2^k} \quad$$
	With $X=\s{a,b}$, we have $\hbd{c}{\s{a,b}} = \kbd{[a,b]}$
\end{example}
\begin{definition}[Limit of a Sequence]
	Let $(X_n)_{n \in \mathbb{N}}$ be a sequence of uncertain numbers.
	We say that the uncertain sequence $(X_n)$ converges to the uncertain number $L$ if:
	For every real number $\epsilon > 0$, there exists a natural number $n_0$ such that for all $n > n_0$, we always have:$$d_H(X_n, L) < \epsilon$$
	In this case, we write \[\lim_{n \to \infty} X_n = L\]
\end{definition}
\begin{definition}[Limit of a Function at a Point]
	Let $f$ be an uncertain function from $\mathcal{X} \rightarrow \mathcal{Y}$ mapping $X \mapsto f(X)$. The uncertain function $f(X)$ is said to have limit $L$ as $X \rightarrow X_0$, denoted as $$\lim\limits_{X \rightarrow X_0}f(X) = L$$
	If $$(\forall \epsilon >0)(\exists \delta >0)(\forall X \in \mathcal{X}) : $$ $$ d_H(X,X_0)<\delta \Rightarrow d_H(f(X),L) < \epsilon$$
\end{definition}
\begin{definition}[Continuous Uncertain Function]
	An uncertain function $f : \mathcal{X} \to \mathcal{Y}$ is called continuous at point $X_0 \in \mathcal{X}$ if $$\lim\limits_{X \rightarrow X_0}f(X) = f(X_0)$$
\end{definition}
\begin{definition}[Difference Quotient]
	Let $f$ be an uncertain function from $\mathcal{X} \rightarrow \mathcal{Y}$ mapping $X \mapsto f(X)$, $h \in \mathbb{K}$.
	The difference quotient of the uncertain function $f(X)$ at $X_0$ is \[\Delta_f(X_0)= \frac{f(X_0 + h) - f(X_0)}{h}\]
\end{definition}
\begin{definition}[Derivative]
	Let $f$ be an uncertain function from $\mathcal{X} \rightarrow \mathcal{Y}$ mapping $X \mapsto f(X)$.
	The derivative of the uncertain function $f(X)$ at $X_0$ is an uncertain number: $$f'(X_0) := \s{\lim_{h\rightarrow 0}{a: a \in \Delta_f(X_0) , \lim_{h\rightarrow 0}a \in \mathbb{K}}}$$
\end{definition}
\begin{example}
	For $f(X) = C \in \tsbd{K}$, we have: $$f'(X) = 0$$
\end{example}
\begin{example}
	For $f(X) = X^2_\times = XX $, we have $f'(X) = X^2_+ = X+X$. Specifically, for instance, taking $X = \s{1,2}$, we have:
	$$f'(X) = X^2_+ = X+X = \s{1,2} + \s{1,2} = \s{2,3,4}$$
\end{example}

\begin{example}
	For $f(X) = (X^2)_1 = X^2 $, we have $f'(X) = 2X$. Specifically, for instance, taking $X = \{1,2\}$, we have:
	$$f'(X) = 2X = 2 \s{1,2} = \s{2,4}$$
\end{example}
\begin{example}
	For $f(X)= (X^2)_1\cdot (X+X)$, we have \[f'(X) = \s{2x_1x_3+2x_2x_3+2x_3x_3 : x_1,x_2,x_3 \in X}\]
	For instance, taking $X=\s{1,2}$, then $f'(X) = \s{6, 8, 10, 16, 20, 24}$.
\end{example}
\begin{theorem}[Derivative of a Sum]
	Let $f : \mathcal{X} \to \mathcal{Y}, g : \mathcal{X} \to \mathcal{Y}$ be uncertain functions differentiable at $X_0 \in \tsbd{K}$. We have \[
		((f+g)'(X_0))_\alpha = (f'(X_0) + g'(X_0))_\alpha, \forall \alpha \in \set{1,m,m'}.
	\]
\end{theorem}
\begin{definition}[Partial Uncertain Integral]
	Let $A,B \in \tsbd{K}$, $f: \mathcal{X} \to \mathcal{Y}$ be a function possessing an antiderivative, and $F$ be an antiderivative of $f$. The partial uncertain integral from $A$ to $B$ of $f(X)$ with respect to antiderivative $F$ is given by:
	\[
		\int_{A}^{B}{f(X)|F} := \s{k = \int_{a}^{b}f(x)dx : a \in F(H_A), b \in F(H_B), k \in \mathbb{K}}
	\]
	Where:
	\begin{align*}
		H_A &:= \left(\bigcup_{x \in A}{F^{-1}(x)}\right)_u \\
		H_B &:= \left(\bigcup_{x \in B}{F^{-1}(x)}\right)_u
	\end{align*}
\end{definition}
\begin{definition}[Uncertain Integral]
	Let $A,B \in \tsbd{K}$, $f: \tsbd{K} \to \tsbd{K}$ be a function possessing antiderivatives. Let $\mathcal{F}$ be the set of all antiderivatives of $f$. Then the uncertain integral from $A$ to $B$ of $f$ is given by \[
	\int_{A}^{B}{f(X)} := \s{\int_{A}^{B}f(X)|F_i : F_i \in \mathcal{F}}
	\]
\end{definition}
\begin{example}
	Evaluate: \[I = \int_{\s{1,2}}^{\s{3,4}}{(X+X)|(X^2)_{m'}}\]
	We have $F(X) = X*X$ as an antiderivative of $f(X)$. We compute:
	\begin{align*}
		F^{-1}(\s{1,2}) &= \s{-1,1,\sqrt{2},-\sqrt{2}} \\
		F^{-1}(\s{3,4}) &= \s{-2,2,\sqrt{3},-\sqrt{3}}
	\end{align*}
	Thus we have: 
	\begin{align*}
		F(F^{-1}(\s{1,2})) &= \s{-1,1,\sqrt{2},-\sqrt{2}} * \s{-1,1,\sqrt{2},-\sqrt{2}} \\
		F(F^{-1}(\s{3,4})) &= \s{-2,2,\sqrt{3},-\sqrt{3}} * \s{-2,2,\sqrt{3},-\sqrt{3}}
	\end{align*}
	Simplifying yields:
	\begin{align*}
		F(F^{-1}(\s{1,2})) &= \s{\pm 1, \pm 2, \pm \sqrt{2}} \\
		F(F^{-1}(\s{3,4})) &= \s{\pm 4, \pm 2\sqrt{3}, \pm 3}
	\end{align*}
	From which we calculate $I$ as equal to:
	$$
		 \s{k = \int_{a}^{b}{2xdx}: a \in \s{\pm 1, \pm 2, \pm \sqrt{2}}, b \in \s{\pm 4, \pm 2\sqrt{3}, \pm 3}, k \in \mathbb{K}}
	$$
	Simplifying yields:
	\[I = \s{5, 7, 8, 10, 11, 12, 14, 15}\]
\end{example}
\begin{theorem}
	Let $A, B \in U(\mathbb{K})$ have canonical forms $A \sim (f_A, d_A)$ and $B \sim (f_B, d_B)$. Let $g(X)$ be an uncertain function with structural antiderivative $F \in \mathcal{F}$ represented by the $k$-variable generating function $f_F(f_{X_1}, \dots, f_{X_k})$.
	
	Consider $k$ differential parameterization paths $f_{X_i}(t_i) = f_A(x_i) + t_i\big(f_B(y_i) - f_A(x_i)\big)$ with $t_i \in [0, 1]$. The total differential form of antiderivative $F$ is:
	\[
	d f_F = \sum_{i=1}^{k} \frac{\partial f_F}{\partial f_{X_i}} \cdot \big(f_B(y_i) - f_A(x_i)\big) dt_i
	\]
	The differential form uncertain integral from $A$ to $B$ of function $g(X)$ corresponding to antiderivative $F$, denoted $\int_A^B g(X)dX \Big| F$, has a canonical form generating function determined by the $k$-fold line integral:
	\[
	f_{I_F}(\mathbf{x}, \mathbf{y}) := \int_{[0, 1]^k} d f_F = f_F\big(f_B(y_1), \dots, f_B(y_k)\big) - f_F\big(f_A(x_1), \dots, f_A(x_k)\big)
	\]
	for all $(\mathbf{x}, \mathbf{y}) \in d_A^k \times d_B^k$.
\end{theorem}

\begin{theorem}
	Let $f,g : \mathcal{X} \to \mathcal{Y}$ be two uncertain functions differentiable at uncertain point $X_0 \in \mathcal{X}$. Suppose representative generating functions of $f(X_0)$ and $g(X_0)$ in canonical form are $f_{f(X_0)(t)}$ and $f_{g(X_0)}(t)$ respectively with $t \in [0,1]^k$.
	\newline
	Then the differential form of the sum function:
	$h(x) = (f+g)(X)$ at $X_0$ is defined by the total linear differential operator:
	\[df_{f+g}(t) = df_{f(X_0)}(t) + df_{g(X_0)}(t)\]
	Explicit representation via partial derivatives over parameter hypercube $[0,1]^k:$
	\[df_{\left(f+g\right)}(t) = \sum_{i=1}^{k}\left[ \left(\frac{\partial f_{f(X_0)}(t)}{\partial t_i} + \frac{\partial f_{g(X_0)}(t)}{\partial t_i}\right) dt_i\right]\]
\end{theorem}